\documentclass[11pt]{article}
\usepackage{fontspec}
\setmainfont{Latin Modern Roman}
\usepackage{amsmath,amssymb,mathtools,unicode-math}
\setmathfont{Latin Modern Math}
\usepackage{xurl,hyperref}
\hypersetup{hidelinks,pdftitle={Problem statement}}
\newcommand{\sref}[2]{\hyperlink{source:#1:#2}{[S#2]}}
\newcommand{\cataloguescope}[1]{\paragraph{Mathematical question.}#1}
\begin{document}
\section{Laplacian Degrees of a Graph}\label{statement}

\subsection{Definitions and mathematical statement}

\paragraph{Shared notation.}$\mathbb N=\{1,2,\ldots\}$, and $\mathbb Z,\mathbb Q,\mathbb R,\mathbb C$ denote the integers, rationals, reals and complex numbers. Any other conventions are specified locally.


For a finite connected simple graph $G$ on $n$ vertices, let $A$ be its adjacency matrix, $D$ the diagonal matrix of degrees, and $L=D-A$ its Laplacian. Order its eigenvalues as $t_1\ge\cdots\ge t_n$, and degrees as $d_1\ge\cdots\ge d_n$. Put $c_k=t_k+k-2$.
\paragraph{Statement recorded in the supplied source.}
The source's recorded conjecture is \[c_k\ge d_k\qquad(1\le k\le n-1).\] Its discussion also asks for a structural characterization of the graphs attaining equality, for each particular index $k$. \sref{statement}{2}
\subsection{Short English statement}
Must every Laplacian degree dominate the corresponding usual degree, and which connected graphs achieve equality?
\subsection{Sources}
\begin{description}
\item[\hypertarget{source:statement:1}{S1}] ULAM.AI and the UnsolvedMath Contributors, UnsolvedMath: A Curated Collection of Open Mathematics Problems, catalogue entry. CC BY 4.0. \url{https://huggingface.co/datasets/ulamai/UnsolvedMath}.
\item[\hypertarget{source:statement:2}{S2}] Open Problem Garden, Laplacian Degrees of a Graph, original node 407, statement and mathematical discussion. \url{https://www.openproblemgarden.org/op/laplacian_degrees_of_a_graph}.
\end{description}
\label{end:statement}
\end{document}
